% topic: latex template for homework

% cmd  : clear; latex fig1.tex  (Note: it will create 3 extra files log,aux and dvi )
% cmd  : clear; xdvi fig1.dvi
% cmd  : clear; dvipdf fig1.dvi

% cmd  : ctrl D if problem occurs
% NOTE : look at end for description of figure in LATEX.
% ref  : https://www.sharelatex.com/learn/Positioning_images_and_tables


\documentclass{article}
\usepackage{blindtext}

%Image-related packages
\usepackage{graphicx}
\usepackage{subcaption}
\usepackage[utf8]{inputenc}
\usepackage[export]{adjustbox}
\usepackage{wrapfig}
%------------------------------

%Table-related commands
\usepackage{array}
\usepackage[table]{xcolor}
\setlength{\arrayrulewidth}{1mm}
\setlength{\tabcolsep}{18pt}
\renewcommand{\arraystretch}{1.5}
\newcolumntype{s}{>{\columncolor[HTML]{AAACED}} p{3cm}}
%-------------------------------------------------------


\title{Positioning Tables and Figures}
\author{ }
\date{\today }

\graphicspath{ {images/} }

\begin{document}
\maketitle
\tableofcontents
\thispagestyle{empty}
\listoffigures
\listoftables
\pagenumbering{arabic}
\clearpage


\section{Positioning images}

This is a sample text.
\clearpage



%%%%%%%%%%%%% graphics environment %%%%%%%%%%%%%%%%%%%%%%

\includegraphics[width=1.0\textwidth, right]{./images/a}


%   we need \usepackage{graphicx} and \usepackage[export]{adjustbox}
%	The package adjustbox enables an additional option in the 
%	\includegraphics command, in the example the picture is 
%	aligned to right. The available values are: 
%	left, right, center, outer and inner, 
%	the last two are intended for two-sided documents. 



%%%%%%%%%%%%%%%%%%%%%%%%%%%%%%%%%%%%%%%%%%%%%%%%%%%%%%%%%%

%%%%%%%%%% figure environment %%%%%%%%%%%%%%%%%
\begin{figure}[ht!]
\includegraphics[width=1.0\textwidth, right]{./images/b.eps}
\caption{fig.2}
\end{figure}
\clearpage



\section{Multiple images in one figure}  % we need \usepackage{subcaption}

	\begin{figure}[h]
 
		\begin{subfigure}{0.5\textwidth}
			\includegraphics[width=0.9\linewidth, height=5cm]{a} 
			\caption{Single Precision}
			\label{fig:subim1}
		\end{subfigure}
		\begin{subfigure}{0.5\textwidth}
			\includegraphics[width=0.9\linewidth, height=5cm]{b}
			\caption{Double precision}
			\label{fig:subim2}
			\end{subfigure}
 
			\caption{Caption for this figure with two images}
			\label{fig:image2}
	\end{figure}.


\section{Wrapping text around a figure}  % needs \usepackage{wrapfig}

	The package wrapfig provides a useful feature, 
	text can be floated around the images.
 
	\begin{wrapfigure}{l}{0.25\textwidth}
		\includegraphics[width=0.9\linewidth]{a} 
		\caption{Single precision}
		\label{fig:subim1}
	\end{wrapfigure}
	
	\begin{verbatim}
	First import the package wrapfig by adding  \usepackage{wrapfig} to the preamble.

	After that you can use the environment <code>wrapfig, 
	it takes two parameters that are passed inside braces: 
	the alignement that can be l, r, c, i or o; 
	this letters stand for left, right, centre, inner and outer 
	(the last two intended for two-sided documents). 
	The second parameter is the width of the figure, in the example 
	is 0.25 the width of the text. See the reference guide for a list of possible length units. 
	\end{verbatim}
\section{Referencing} \label{myref}
  We can reference images, for instance, the image \ref{fig:subim1} shows single precision graph.
  \clearpage
  
\section{Math references} \label{mathrefs}
As mentioned in section \ref{myref}, different elements can 
be referenced within a document
 
	\subsection{powers series} \label{subsection}
 
		\begin{equation} \label{eq:1}
			\sum_{i=0}^{\infty} a_i x^i
		\end{equation}
	
	\subsection{example of reference}
	The equation (\ref{eq:1}) is a typical power series.


 

\end{document}

%%%%%%%%%%%%%%%%%%%%%%%%%%%%%%%%%%%%%%%%%%%%%%%%%%%%%%%%%%%%%%
%\includegraphics[attr1=val1, attr2=val2, ..., attrn=valn]{imagename}
%\DeclareGraphicsExtensions{.pdf,.png,.jpg}
%
%\includegraphics[width=\linewidth]{chick}
%\includegraphics[width=\textwidth]{chick}
%\includegraphics[height=\textheight]{chick}
%\includegraphics[scale=0.5, angle=180]{chick}
%
%
%% for borders
%\setlength\fboxsep{0pt}
%\setlength\fboxrule{0.5pt}
%\fbox{\includegraphics{chick}}
























